\documentclass{jsarticle}
\usepackage[dvipdfmx,final]{graphicx}
\usepackage{amsmath}
\usepackage{amssymb}
\begin{document}
\title{Topology theorem from string theory and \\
       dualty element of differential structure}
\author{Masaaki Yamaguchi}
\date{}
\maketitle


  $$ \int d \tau = f(x), \tau = (\alpha \circ \beta)^{-1}\cdot \alpha \beta $$
  $$ \ge {{\alpha \beta} \over {\alpha^{-1}\beta^{-1}}} $$
  $$ (\alpha \circ \beta)^{-1} \le \alpha^{-1} + \beta^{-1} $$
  $$ \beta(p,q) = {\Gamma(p) \Gamma(q) \over \Gamma(p+q)} $$
  $$ {{\int \int d \tau} \over {\int d \tau}} = \begin{vmatrix} {{\int \int d 
  \tau} \over {\int d \tau}} \end{vmatrix}d \psi = \begin{pmatrix} {{
  \nabla_i \nabla_j f} \over {\Delta f}} \end{pmatrix}^{-1} $$
  $$ (f(x)g(x))^{'} = f^{'}(x)g(x) + f(x)g^{'}(x) $$
  $$ \int {\begin{pmatrix} f(x)g(x) \end{pmatrix}}^{'} = 
	  \int f^{'}(x)g(x) + \int f(x) g^{'}(x) = 
	  \int f(x)g(x) = [f(x)g(x)] - \int f^{'}(x)g(x) $$
	  $$ [f(x),g(x)] = f(x)g(x) - g(x)f(x) $$
	  $$ (fg)^{'} = fg^{'} + g^{'}f, \begin{pmatrix}{f \over g}
	  \end{pmatrix}^{'} = {{f^{'}g - g^{'}f} \over {g^2}}$$
	  $$ {{\{f,g\}} \over {[f,g]}} = {{(fg)^{'}\otimes dx_{fg}}
	  \over {{({f \over g})^{'}} \otimes g^{-2}dx_{fg}}} $$
	  $$ \begin{pmatrix} {{\{f,g\}} \over {[f,g]}}\end{pmatrix} = (i)^{'} $$
	これらの方程式は、フロベニウスの定理と実数条件からの
		  大域的微分方程式の最小単位による積分でもある。
		  これらは、演算子同士による計算、記号論に行き着く。
		  非可換代数方程式、Lie代数の部分群、
		  フェルミオンの部分群、
		  どこから始めたらいいのか。神の人刺し、神の見えざる手、
		  フェルミオンとボソンの商代数による積と商の微分部分群
		  による準同型写像の計算についての幾通りの算術式。
		   
		  $$ V(M) = {\partial \over \partial f}
		  \begin{pmatrix}{}^{N}\int [f \setminus M]^{\oplus N}
		  \end{pmatrix}^{\mu\nu}dx^{\mu}dx^{\nu} $$
		  $$ \bigoplus_{k=0}^{\infty} [f \setminus g] = \vee(M
		  \wedge N) $$
		  $$ [f \setminus g] = (F - g) \otimes {1 \over M} $$
		  $$ V - E + F = 2, \log {\sin \theta \over i\cos \theta}
		  = \log(\sin \theta) - \log(i \cos \theta) $$
		  $$ = - 2 R_{ij}, {1 \over i}\tan \theta = H \psi  $$
		  $$ e^{i \theta} = \cos \theta + i \sin \theta $$
		  $$ \vee \int{{C^+ \nabla M_m} \over {\Delta(M^+_- \nabla
		  C^+_-)}}i = \exists(M^+_- \nabla R^+) $$
		  $$ \vee(R + \nabla_i \nabla_j f)^n = \int{{\wedge(R + 
		  \nabla_i \nabla_j f)^2} \over 
		 {\exists(R + \Delta f)^n}} $$
		 $$ \wedge(R + \nabla_i \nabla_j f)^2 = {d \over df}
		 {\int \int {1 \over (y \log y)^{1 \over 2}}dy_m} $$
		 $$ {d \over dt}g_{ij}(t) = - 2 R_{ij} $$
		 $$ \vee \int \wedge (R + \nabla_i \nabla_j f)^x
		 = {{\wedge(R + \nabla_i \nabla_j f)^n }\over 
		 {\exists(R + \nabla_i \nabla_j f \circ g)}}$$
		  $$\bigoplus f_n^m = f_1 \oplus f_2 \oplus f_3 \oplus \cdots
		  \bigoplus f_n^m $$
		  $$ \cong \wedge[f\oplus g] \cong \wedge[f \bigoplus M^+_1] $$
		  $$ \cong \wedge[f \bigoplus M^+_-], \sum_{k=0}^{\infty}
		  {1 \over n^s} $$

\includegraphics[width=15cm,clip]{string_equation2}
\includegraphics[width=15cm,clip]{string_theorem}
\includegraphics[width=15cm,clip]{world_theorem}
\includegraphics[width=15cm,clip]{gamma_theorem}
\includegraphics[width=15cm,clip]{Knot_theorem_topology}




\end{document}
